\section{Hardware-Aware Implementations}

To achieve real-time performance on mobile and edge devices, the proposed system incorporates hardware-aware optimizations tailored for both CPU and GPU execution. While the underlying algebraic formulation remains consistent across implementations, computational stages are reorganized to exploit the architectural strengths—and mitigate the bottlenecks—of each processing unit.

\subsection{CPU Implementation}

\begin{figure}[htbp]
\centering
\resizebox{\textwidth}{!}{%
\begin{tikzpicture}[x=1cm, y=1cm]

  %% ══════════════════════════════════════════════════════════════
  %%  TOP PIPELINE (3 Vertical Rows)
  %% ══════════════════════════════════════════════════════════════
  
  \begin{scope}[node distance=1.0cm]
      %% ROW 2: The 4 SIMD-accelerated steps (Placed first for central anchoring)
      \node[simdhw] (GS) at (0,0) {Grayscale\\[-2pt]\tiny\itshape NEON/SSE2};
      \node[simdhw, right=of GS] (KP) {FAST Keypoints\\[-2pt]\tiny\itshape NEON/SSE4};
      \node[simdhw, right=of KP] (DESC) {Binary Descriptors\\[-2pt]\tiny\itshape NEON/AVX2};
      \node[simdnode, right=of DESC] (MATCH) {Hamming Match\\[-2pt]\tiny\itshape popcount/VCNT};

      %% ROW 1: Camera Frame (Aligned directly above Grayscale)
      \node[bbnodeD, above=2.0cm of GS] (IN) {Camera Frame};

      %% ROW 3: Homography (Aligned directly below Match)
      \node[bbnodeD, below=2.0cm of MATCH] (POSE) {Homography /\\[-2pt]Optical Flow};
  \end{scope}

  %% Connections
  % Downward arrow into pipeline
  \draw[fwB, line width=1.2pt, color=bbdark] (IN.south) -- (GS.north);
      
  % Horizontal SIMD internal arrows
  \draw[fwB] (GS.east) -- (KP.west);
  \draw[fwB] (KP.east) -- (DESC.west);
  \draw[fwB] (DESC.east) -- (MATCH.west);
  
  % Downward arrow out of pipeline
  \draw[fwB, line width=1.2pt, color=bbdark] (MATCH.south) -- (POSE.north);

  %% Brace below pipeline
  \draw[decorate, decoration={brace,amplitude=5pt,mirror}, simdborder, line width=0.7pt]
      ($(GS.south west)+(0,-0.2)$) -- ($(MATCH.south east)+(0,-0.2)$)
      node[midway, below=6pt, font=\scriptsize\sffamily, text=simdborder]
          {SIMD-accelerated stages \ (OpenCV NEON / SSE / AVX intrinsics)};

  %% ══════════════════════════════════════════════════════════════
  %%  BOTTOM PANEL (Distributed to Full Width)
  %% ══════════════════════════════════════════════════════════════
  
  \begin{scope}[shift={($(GS.south west) + (0, -4.5cm)$)}]
      
      % Panel Title
      \node[font=\footnotesize\sffamily\bfseries, text=bbdark!80, anchor=north west] 
          (ILL_TITLE) at (-0.5, 0.5) {How SIMD parallelises pixel processing};

      %% --- SECTION 1: SCALAR PANEL (Left) ---
      \begin{scope}[shift={(0, -1)}]
          \node[font=\tiny\sffamily\bfseries, text=bbdark, anchor=south west] 
              at (0, 0.25) {Scalar --- 1 pixel / cycle};
          
          \foreach \row/\actcol in {0/0, 1/1, 2/2, 3/3} {
              \foreach \col in {0,...,7} {
                  \ifnum\col=\actcol
                      \fill[bbdark!75] (\col*0.48, -\row*0.38) rectangle +(0.44,0.30);
                  \else
                      \fill[scalarclr,draw=bbborder!50,line width=0.2pt] (\col*0.48, -\row*0.38) rectangle +(0.44,0.30);
                  \fi
              }
              \node[font=\tiny\sffamily,text=bbdark,anchor=west] 
                  at (8*0.48, -\row*0.38 + 0.15) {cycle \number\numexpr\row+1\relax};
          }
          \node[font=\tiny\sffamily\itshape,text=bbdark!65,anchor=north] 
              at (1.8, -1.15) {\ldots\ 8 cycles total for 8 pixels};
      \end{scope}

      %% --- ARROW 1 (Vectorise) ---
      \draw[-{Stealth[length=8pt,width=7pt]}, line width=1.4pt, simdborder]
          (5.0, -1.6) -- (6.8, -1.6)
          node[midway, above, font=\tiny\sffamily\bfseries, text=simdborder] {vectorise};

      %% --- SECTION 2: VECTOR PANEL (Center) ---
      \begin{scope}[shift={(7.0, -1)}]
          \node[font=\tiny\sffamily\bfseries, text=simdborder, anchor=south west] 
              at (0, 0.3) {SIMD Vector --- 8 pixels / cycle};
          
          \foreach \col in {0,...,7} {
              \fill[vectorclr,draw=simdborder!50,line width=0.3pt] (\col*0.48, 0) rectangle +(0.44,0.30);
              \node[font=\tiny\sffamily\bfseries,text=simdborder!85!black] 
                  at (\col*0.48 + 0.22, 0.15) {\number\numexpr\col+1\relax};
          }
          \node[font=\tiny\sffamily,text=simdborder,anchor=west] 
              at (8*0.48, 0.15) {cycle 1 only!};

          %% Register outline
          \draw[simdborder!70, line width=0.6pt, rounded corners=2pt]
              (-0.05, -0.05) rectangle (8*0.48-0.05, 0.35);
          \node[font=\tiny\sffamily\bfseries,text=simdborder!75,anchor=north]
              at (1.9, -0.1) {128-bit NEON / 256-bit AVX register};

          %% Speedup badge
          \node[draw=simdborder, fill=simdblue!20, rounded corners=4pt,
                font=\normalsize\sffamily, text=simdborder, inner sep=5pt]
              at (1.9, -0.8) {$\times 4$--$8\times$ faster};

          \node[font=\tiny\sffamily\itshape, text=simdborder!70, anchor=north]
              at (1.9, -1.2) {all 8 pixels processed in a single instruction};
      \end{scope}

      %% --- ARROW 2 (Dispatch) ---
      \draw[-{Stealth[length=8pt,width=7pt]}, line width=1.4pt, simdborder]
          (11.0, -1.6) -- (13.8, -1.6)
          node[midway, above, font=\tiny\sffamily\bfseries, text=simdborder] {executed on};

      %% --- SECTION 3: CPU CHIP (Right) ---
      \begin{scope}[shift={(14.0, -1.5)}]
          \fill[simdblue!30, draw=simdborder, line width=1.0pt, rounded corners=4pt]
              (0, 0) rectangle (1.8, 1.2);

          %% CPU Cores
          \foreach \x in {0.2, 0.7, 1.2} {
              \foreach \y in {0.15, 0.5, 0.85} {
                  \fill[vectorclr,draw=simdborder!40,line width=0.25pt] 
                      (\x, \y) rectangle +(0.4,0.25);
              }
          }

          %% Pins top/bottom
          \foreach \px in {0.2, 0.48, 0.76, 1.04, 1.32, 1.6} {
              \draw[simdborder,line width=0.6pt] (\px, 1.2) -- (\px, 1.35);
              \draw[simdborder,line width=0.6pt] (\px, 0) -- (\px, -0.15);
          }
          %% Pins sides
          \foreach \py in {0.15, 0.4, 0.65, 0.9} {
              \draw[simdborder,line width=0.6pt] (0, \py) -- (-0.15, \py);
              \draw[simdborder,line width=0.6pt] (1.8, \py) -- (1.95, \py);
          }

          \node[font=\scriptsize\sffamily\bfseries, text=simdborder]
              at (0.9, -0.4) {ARM / x86 CPU};
          \node[font=\tiny\sffamily\bfseries, text=simdborder!70]
              at (0.9, -0.65) {NEON / SSE / AVX units};
              
          \node (CPU_ANCHOR) at (1.95, -0.65) {};
      \end{scope}

      %% Invisible bounding box anchors perfectly aligned with the top pipeline width
      \coordinate (PANEL_LEFT) at (GS.west |- ILL_TITLE.north);
      \coordinate (PANEL_RIGHT) at (MATCH.east |- CPU_ANCHOR);
      \node (PANEL_BBOX) [fit=(PANEL_LEFT) (PANEL_RIGHT)] {};

  \end{scope}

  %% ══════════════════════════════════════════════════════════════
  %%  BACKGROUND PANELS (Automatically sized)
  %% ══════════════════════════════════════════════════════════════
  \begin{pgfonlayer}{background}
    %% Pipeline outer strip
    \node[fill=bbgray!20, draw=bbborder!35, rounded corners=7pt,
          inner sep=12pt, fit=(IN)(GS)(KP)(DESC)(MATCH)(POSE)] (pipelinepanel) {};

    %% SIMD blue highlight over accelerated nodes
    \node[fill=simdblue!20, draw=simdborder!30, rounded corners=6pt,
          inner sep=9pt, fit=(GS)(KP)(DESC)(MATCH)] {};

    %% Bottom Illustration panel
    \node[fill=bbgray!20, draw=bbborder!25, rounded corners=6pt,
          inner xsep=10pt, inner ysep=10pt, fit=(PANEL_BBOX)] {};
  \end{pgfonlayer}

\end{tikzpicture}%
}
\caption{CPU/SIMD pipeline configuration. Grayscale conversion, FAST keypoint
detection, binary descriptor computation, and Hamming-distance matching are
vectorised via SIMD intrinsics (ARM NEON on mobile, SSE4/AVX2 on desktop).
The lower panel contrasts scalar execution---one pixel per cycle across 8
cycles---with SIMD vector execution where a single 128-bit NEON or 256-bit
AVX2 instruction processes all 8 pixels simultaneously, yielding
$4$--$8\times$ throughput on these stages.}
\label{fig:cpu}
\end{figure}

The CPU-based pipeline is designed to minimize memory latency and maximize instruction-level parallelism. Image frames are acquired directly from the native camera interface, enabling operation on CPU-resident buffers and avoiding costly data transfers. This zero-copy design eliminates synchronization overhead associated with GPU offloading.

As illustrated in Figure~\ref{fig:cpu}, the most computationally intensive stages—grayscale conversion, FAST keypoint detection, binary descriptor construction, and Hamming-distance matching—are accelerated using Single Instruction, Multiple Data (SIMD) vectorization. By leveraging 128-bit ARM NEON intrinsics on mobile processors (or 256-bit AVX2 on desktop architectures), the pipeline transitions from scalar processing to vectorized execution, enabling multiple pixels to be processed per cycle.

Binary descriptor matching is particularly well-suited to SIMD execution. Descriptor comparisons are implemented using vectorized bitwise XOR operations followed by hardware population-count instructions (e.g., \texttt{VCNT} or \texttt{popcount}). Combined with cache-aligned memory layouts and compiler-assisted auto-vectorization, this approach yields a $4\times$–$8\times$ improvement in feature processing throughput.

Once correspondences are established, motion estimation—via RANSAC-based homography estimation and pyramidal Lucas--Kanade optical flow—is executed on the CPU. These stages involve iterative and branch-heavy computations, which benefit from the CPU's low-latency cache hierarchy and efficient branch prediction.

\subsection{GPU Implementation}

\begin{figure}[htbp]
\centering
\resizebox{\textwidth}{!}{%
\begin{tikzpicture}[x=1cm,y=1cm]

  %% ══════════════════════════════════════════════════════════════
  %%  TOP PIPELINE (3 Vertical Rows)
  %% ══════════════════════════════════════════════════════════════
  
  \begin{scope}[node distance=1.0cm]
      %% ROW 2: The 4 GPU-accelerated steps (Placed first for central anchoring)
      \node[gpuhw] (GS3) at (0,0) {Grayscale\\[-2pt]\scriptsize\itshape GPU shader};
      \node[gpuhw, right=of GS3] (KP3) {Keypoint Detection\\[-2pt]\scriptsize\itshape shader kernel};
      \node[gpuhw, right=of KP3] (DESC3) {Descriptor Compute\\[-2pt]\scriptsize\itshape parallel threads};
      \node[gpuhw, right=of DESC3] (MATCH3) {Parallel Match\\[-2pt]\scriptsize\itshape thread-per-descriptor};

      %% ROW 1: Camera Frame (Aligned above Grayscale)
      \node[bbnodeD, above=2.0cm of GS3] (IN3) {Camera Frame};

      %% ROW 3: Homography (Aligned below Match)
      \node[bbnodeD, below=2.0cm of MATCH3] (POSE3) {Homography/\\[-2pt]Optical Flow};
  \end{scope}

  %% Connections
  % Downward arrow into GPU
  \draw[fwG, line width=1.2pt] (IN3.south) -- (GS3.north)
      node[annG, midway, right, align=left, xshift=2pt] {\textbf{PCIe Upload}\\Frame to VRAM};
      
  % Horizontal GPU internal arrows
  \draw[fwG] (GS3.east) -- (KP3.west);
  \draw[fwG] (KP3.east) -- (DESC3.west);
  \draw[fwG] (DESC3.east) -- (MATCH3.west);
  
  % Downward arrow out of GPU
  \draw[fwG, line width=1.2pt, color=bbdark] (MATCH3.south) -- (POSE3.north)
      node[annG, midway, left, align=right, xshift=-2pt, color=bbdark] {\textbf{PCIe Readback}\\Matches to RAM};

  %% Brace below GPU pipeline
  \draw[decorate, decoration={brace,amplitude=5pt,mirror}, gpuborder, line width=0.7pt]
      ($(GS3.south west)+(0,-0.2)$) -- ($(MATCH3.south east)+(0,-0.2)$)
      node[midway, below=6pt, font=\scriptsize\sffamily, text=gpuborder]
          {GPU-accelerated stages (on-device GLSL/Metal shaders)};

  %% ══════════════════════════════════════════════════════════════
  %%  BOTTOM PANEL (Distributed to Full Width)
  %% ══════════════════════════════════════════════════════════════
  
  \begin{scope}[shift={($(GS3.south west) + (0, -5cm)$)}]
      
      % Panel Title
      \node[font=\footnotesize\sffamily\bfseries, text=bbdark!80, anchor=north west] 
          (GPU_TITLE) at (0, 1) {How GPU offloading parallelises feature tracking};

      %% --- SECTION 1: MEMORY TRANSFER (Left) ---
      \begin{scope}[shift={(0, -1)}]
          \node[font=\tiny\sffamily\bfseries, text=bbdark, anchor=south] 
              at (1.9, 0.9) {1. Memory Transfer Bottleneck};
              
          % CPU RAM Box
          \draw[bbborder, fill=bbgray!60, rounded corners=2pt] (0, 0) rectangle (1.1, 0.8);
          \node[font=\tiny\sffamily, text=bbdark, align=center] at (0.55, 0.4) {CPU\\RAM};
          
          % GPU VRAM Box
          \draw[gpuborder!70, fill=gpugreen!40, rounded corners=2pt] (2.7, 0) rectangle (3.8, 0.8);
          \node[font=\tiny\sffamily, text=gpuborder!60!black, align=center] at (3.25, 0.4) {GPU\\VRAM};
          
          % Arrows
          \draw[-{Stealth[length=6pt,width=5pt]}, line width=1.2pt, gpuborder] 
              (1.2, 0.55) -- (2.6, 0.55) node[midway, above, font=\tiny\sffamily\itshape] {Upload};
          \draw[{Stealth[length=6pt,width=5pt]}-, line width=1.2pt, bbdark!60] 
              (1.2, 0.25) -- (2.6, 0.25) node[midway, below, font=\tiny\sffamily\itshape] {Readback};
      \end{scope}

      %% --- SECTION 2: HARDWARE / SHADER CORES (Center) ---
      \begin{scope}[shift={(6.5, -0.6)}]
          \node[font=\tiny\sffamily\bfseries, text=gpuborder, anchor=south] 
              at (2.5, 0.6) {2. Shader Multiprocessors (Massive Grid)};
          
          %% GPU Chip illustration
          \coordinate (gchip) at (0, -0.2);
          \fill[gpugreen!30, draw=gpuborder, line width=0.8pt, rounded corners=3pt]
              ($(gchip)+(-0.7,-0.6)$) rectangle ($(gchip)+(0.7,0.6)$);
          \foreach \p in {-0.42, 0, 0.42}{
            \draw[gpuborder, line width=0.6pt] ($(gchip)+(\p,0.6)$) -- ($(gchip)+(\p,0.8)$);
            \draw[gpuborder, line width=0.6pt] ($(gchip)+(\p,-0.6)$) -- ($(gchip)+(\p,-0.8)$);
          }
          \foreach \p in {-0.3, 0.3}{
            \draw[gpuborder, line width=0.6pt] ($(gchip)+(-0.7,\p)$) -- ($(gchip)+(-0.9,\p)$);
            \draw[gpuborder, line width=0.6pt] ($(gchip)+(0.7,\p)$) -- ($(gchip)+(0.9,\p)$);
          }
          \node[font=\scriptsize\sffamily\bfseries, text=gpuborder!80!black] at (gchip) {GPU};
          
          %% Core Grid (representing thousands of ALUs)
          \begin{scope}[shift={(1.4, -0.8)}]
              \foreach \row in {0,1,2,3}{
                \foreach \col in {0,...,8}{
                  \fill[gpucore, opacity=0.85] (\col*0.42, \row*0.35) rectangle +(0.36,0.28);
                  \draw[gpuborder!40, line width=0.25pt] (\col*0.42, \row*0.35) rectangle +(0.36,0.28);
                }
              }
          \end{scope}
      \end{scope}

      %% --- SECTION 3: THREAD EXECUTION (Right) ---
      \begin{scope}[shift={(13.2, -1)}]
          \node[font=\tiny\sffamily\bfseries, text=gpuborder, anchor=south] 
              at (1.2, 1) {3. Thread-per-Descriptor Mapping};
              
          % Input array
          \foreach \x in {0,1,2,3,4} {
              \draw[bbborder!70, fill=bbgray!80, rounded corners=1pt] (\x*0.6, 0.6) rectangle +(0.5, 0.3);
          }
          \node[font=\tiny\sffamily, text=bbdark, anchor=east] at (-0.1, 0.75) {Input Array};
          
          % Threads
          \foreach \x in {0,1,2,3,4} {
              \draw[-{Stealth[length=4pt,width=3pt]}, gpuborder, line width=0.8pt] 
                  (\x*0.6 + 0.25, 0.5) -- (\x*0.6 + 0.25, 0.1);
              \node[font=\tiny\sffamily\bfseries, text=gpuborder] at (\x*0.6 + 0.25, 0.3) {$\sim$}; 
          }
          \node[font=\tiny\sffamily, text=gpuborder, anchor=east] at (-0.1, 0.3) {GPU Threads};
          
          % Output array
          \foreach \x in {0,1,2,3,4} {
              \draw[gpuborder!50, fill=gpugreen!40, rounded corners=1pt] (\x*0.6, -0.3) rectangle +(0.5, 0.3);
          }
          \node[font=\tiny\sffamily, text=gpuborder!60!black, anchor=east] at (-0.1, -0.15) {Results};
          
          \node[font=\tiny\sffamily\itshape, text=bbdark!70, anchor=north] 
              at (1.1, -0.45) {Thousands of threads execute in unison};
              
          % Invisible anchor for bounding box
          \node (THREAD_ANCHOR) at (3.5, -0.45) {};
      \end{scope}

      %% Invisible bounding box perfectly aligned with the top pipeline width
      \coordinate (PANEL_LEFT) at (IN3.west |- GPU_TITLE.north);
      \coordinate (PANEL_RIGHT) at (POSE3.east |- THREAD_ANCHOR);
      \node (GPU_PANEL_BBOX) [fit=(PANEL_LEFT) (PANEL_RIGHT)] {};

  \end{scope}

  %% ══════════════════════════════════════════════════════════════
  %%  BACKGROUND PANELS
  %% ══════════════════════════════════════════════════════════════
  \begin{pgfonlayer}{background}
    %% Pipeline outer strip
    \node[fill=bbgray!20, draw=bbborder!35, rounded corners=7pt,
          inner sep=12pt, fit=(IN3)(GS3)(KP3)(DESC3)(MATCH3)(POSE3)] (pipelinepanel) {};

    %% GPU green highlight over accelerated nodes
    \node[fill=gpugreen!15, draw=gpuborder!35, rounded corners=6pt,
          inner sep=9pt, fit=(GS3)(KP3)(DESC3)(MATCH3)] {};

    %% Bottom Illustration panel
    \node[fill=bbgray!20, draw=bbborder!25, rounded corners=6pt,
          inner xsep=10pt, inner ysep=10pt, fit=(GPU_PANEL_BBOX)] {};
  \end{pgfonlayer}

\end{tikzpicture}%
}
\caption{GPU-accelerated configuration. The pipeline is restructured so that Grayscale conversion, Keypoint detection, descriptor computation, and Hamming-distance matching are all offloaded to programmable GPU shaders. The camera frame is uploaded once to VRAM. Thousands of threads process the image in parallel, and only the compact match results are read back to the CPU for homography and optical flow.}
\label{fig:gpu}
\end{figure}

The GPU implementation offloads feature extraction and descriptor matching to programmable shader pipelines, trading low-latency sequential execution for high-throughput parallel processing.

As shown in Figure~\ref{fig:gpu}, a key challenge in GPU acceleration is minimizing memory transfer overhead. To address this, the camera frame is uploaded to GPU memory once per frame, after which all intermediate processing is performed entirely on the GPU using a sequence of GLSL or Metal shaders with ping-pong buffering. This design ensures that high-bandwidth intermediate data remains in GPU memory.

The GPU pipeline parallelizes the following stages:

\begin{itemize}
    \item \textbf{Grayscale Conversion and Smoothing:} Implemented as a fragment shader pass.
    \item \textbf{Keypoint Detection and NMS:} Parallel evaluation of local intensity contrasts for FAST corner detection.
    \item \textbf{Descriptor Extraction:} Efficient texture sampling to construct binary descriptors.
    \item \textbf{Parallel Matching:} A compute shader employing a thread-per-descriptor model, where thousands of threads simultaneously compute Hamming distances across descriptor sets.
\end{itemize}

While GPUs excel at data-parallel operations, iterative algorithms such as RANSAC and pyramidal Lucas--Kanade tracking suffer from control-flow divergence and synchronization overhead. Consequently, the system adopts a hybrid execution strategy: the GPU reduces dense image data into a compact set of correspondences, which are then transferred back to the CPU. The CPU subsequently performs motion estimation using its efficient control-flow handling.

Together, these hardware-aware implementations enable the system to adapt to the available compute resources, balancing SIMD-optimized CPU execution with massively parallel GPU processing. This hybrid design ensures efficient feature processing while maintaining real-time performance across a range of mobile and embedded platforms.